% GENERATED FILE. Edit canonical lessons, then run npm run generate:books.
% canonical-source-hash: fnv1a64:ddb1fc8a5381c7c1
% canonical-lessons: GE-C429-schimmel, GE-C429-flackern, GE-C429-klemmen, GE-C429-quietschen, GE-C429-defekt, GE-R429-review-422-429

\chapter{Kleine Schäden}
\label{ch:ge-kleine-schaden}
\hlchaptermodality{429}

\begin{chapteropening}
\textbf{By the end of this chapter:} \emph{I can describe small problems in a room: mould, a flickering light, a door that sticks, a squeak and something broken.}

Complete the last lesson of chapter 429: I can describe small problems in a room: mould, a flickering light, a door that sticks, a squeak and something broken.
\end{chapteropening}

\section[der Schimmel]{der Schimmel — the mould (US: mold); a white horse}

\label{lesson:GE-C429-schimmel}



\begin{quote}
Before the new one: say the German for leaky, then the German for the plumber.
\end{quote}

\subsection*{You'll want to know: der Schimmel}
\textbf{der Schimmel} — \textquotedblleft{}the mould (US: mold); a white horse\textquotedblright{}.

\begin{quote}
In der Ecke vom Bad ist Schimmel.
\end{quote}

\begin{quote}
\emph{There is mould in the corner of the bathroom.}
\end{quote}

No plural for mould. \emph{Ein Schimmel} is also a white horse.

\begin{cousinweb}
\emph{Schimmel} is inherited from Middle High German \emph{schimel} (EtymDB). Dutch \emph{schimmel} 'mould' is a cousin; the white horse is usually said to be named for its mould-white colour.
\end{cousinweb}

\begin{grammarlens}[title={What you've built}]
One more word for telling, explaining and describing at B1.
\end{grammarlens}

\subsection*{Your turn}
\begin{itemize}
\raggedright
  \item \emph{Say it:} \emph{der Schimmel}
  \item \emph{Say it:} \emph{der Schimmel}, once more
  \item \emph{Say it:} a sentence of your own with \emph{der Schimmel}
  \item \emph{Recall:} say the German for leaky, then the German for the plumber, then say \emph{der Schimmel} again
\end{itemize}

\begin{tcolorbox}[breakable,colback=teal!4,colframe=teal!35!black,title={Before you move on}]
What does der Schimmel mean? (\textquotedblleft{}The mould (US: mold); a white horse\textquotedblright{}.) Say it once more.
\end{tcolorbox}
\section[flackern]{flackern — to flicker}

\label{lesson:GE-C429-flackern}



\begin{quote}
Before the new one: say the German for the plumber, then the German for the mould.
\end{quote}

\subsection*{You'll want to know: flackern}
\textbf{flackern} — \textquotedblleft{}to flicker\textquotedblright{}.

\begin{quote}
Die Lampe über dem Bett flackert.
\end{quote}

\begin{quote}
\emph{The lamp above the bed flickers.}
\end{quote}

\emph{Es flackert}, \emph{es hat geflackert}.

\begin{cousinweb}
\emph{flackern} is an old German verb for a light that flutters (EtymDB lists 'to flicker, to flutter'); its origin before that is not settled.
\end{cousinweb}

\begin{grammarlens}[title={What you've built}]
One more word for telling, explaining and describing at B1.
\end{grammarlens}

\subsection*{Your turn}
\begin{itemize}
\raggedright
  \item \emph{Say it:} \emph{flackern}
  \item \emph{Say it:} \emph{flackern}, once more
  \item \emph{Say it:} a sentence of your own with \emph{flackern}
  \item \emph{Recall:} say the German for the plumber, then the German for the mould, then say \emph{flackern} again
\end{itemize}

\begin{tcolorbox}[breakable,colback=teal!4,colframe=teal!35!black,title={Before you move on}]
What does flackern mean? (\textquotedblleft{}To flicker\textquotedblright{}.) Say it once more.
\end{tcolorbox}
\section[klemmen]{klemmen — to jam, to stick; to pinch}

\label{lesson:GE-C429-klemmen}



\begin{quote}
Before the new one: say the German for the mould, then the German for to flicker.
\end{quote}

\subsection*{You'll want to know: klemmen}
\textbf{klemmen} — \textquotedblleft{}to jam, to stick; to pinch\textquotedblright{}.

\begin{quote}
Die Tür zum Balkon klemmt.
\end{quote}

\begin{quote}
\emph{The door to the balcony sticks.}
\end{quote}

\emph{Ich habe mir den Finger geklemmt} — I caught my finger.

\begin{cousinweb}
\emph{klemmen} is an inherited Germanic verb for squeezing and pinching (EtymDB lists 'to be stuck, to stick'). English \emph{clamp} is usually counted among its cousins.
\end{cousinweb}

\begin{grammarlens}[title={What you've built}]
One more word for telling, explaining and describing at B1.
\end{grammarlens}

\subsection*{Your turn}
\begin{itemize}
\raggedright
  \item \emph{Say it:} \emph{klemmen}
  \item \emph{Say it:} \emph{klemmen}, once more
  \item \emph{Say it:} a sentence of your own with \emph{klemmen}
  \item \emph{Recall:} say the German for the mould, then the German for to flicker, then say \emph{klemmen} again
\end{itemize}

\begin{tcolorbox}[breakable,colback=teal!4,colframe=teal!35!black,title={Before you move on}]
What does klemmen mean? (\textquotedblleft{}To jam, to stick; to pinch\textquotedblright{}.) Say it once more.
\end{tcolorbox}
\section[quietschen]{quietschen — to squeak, to creak; to screech}

\label{lesson:GE-C429-quietschen}



\begin{quote}
Before the new one: say the German for to flicker, then the German for to jam.
\end{quote}

\subsection*{You'll want to know: quietschen}
\textbf{quietschen} — \textquotedblleft{}to squeak, to creak; to screech\textquotedblright{}.

\begin{quote}
Das alte Bett quietscht bei jeder Bewegung.
\end{quote}

\begin{quote}
\emph{The old bed squeaks every time you move.}
\end{quote}

\emph{Die Reifen quietschen} — the tyres screech.

\begin{cousinweb}
\emph{quietschen} is a sound word that imitates a high, sharp squeak. EtymDB does not list it.
\end{cousinweb}

\begin{grammarlens}[title={What you've built}]
One more word for telling, explaining and describing at B1.
\end{grammarlens}

\subsection*{Your turn}
\begin{itemize}
\raggedright
  \item \emph{Say it:} \emph{quietschen}
  \item \emph{Say it:} \emph{quietschen}, once more
  \item \emph{Say it:} a sentence of your own with \emph{quietschen}
  \item \emph{Recall:} say the German for to flicker, then the German for to jam, then say \emph{quietschen} again
\end{itemize}

\begin{tcolorbox}[breakable,colback=teal!4,colframe=teal!35!black,title={Before you move on}]
What does quietschen mean? (\textquotedblleft{}To squeak, to creak; to screech\textquotedblright{}.) Say it once more.
\end{tcolorbox}
\section[defekt]{defekt — faulty, not working (of a device)}

\label{lesson:GE-C429-defekt}



\begin{quote}
Before the new one: say the German for to jam, then the German for to squeak.
\end{quote}

\subsection*{You'll want to know: defekt}
\textbf{defekt} — \textquotedblleft{}faulty, not working (of a device)\textquotedblright{}.

\begin{quote}
Der Fernseher im Zimmer ist defekt.
\end{quote}

\begin{quote}
\emph{The television in the room is not working.}
\end{quote}

A little formal; everyday speech says \emph{kaputt}. \emph{Der Defekt} — the fault.

\begin{cousinweb}
\emph{defekt} is a learned borrowing from Latin \emph{dēfectus} 'weakened, lacking', the participle of \emph{dēficere} 'to fail, to fall short' (Lewis \& Short). EtymDB does not list it.
\end{cousinweb}

\begin{grammarlens}[title={What you've built}]
One more word for telling, explaining and describing at B1.
\end{grammarlens}

\subsection*{Your turn}
\begin{itemize}
\raggedright
  \item \emph{Say it:} \emph{defekt}
  \item \emph{Say it:} \emph{defekt}, once more
  \item \emph{Say it:} a sentence of your own with \emph{defekt}
  \item \emph{Recall:} say the German for to jam, then the German for to squeak, then say \emph{defekt} again
\end{itemize}

\begin{tcolorbox}[breakable,colback=teal!4,colframe=teal!35!black,title={Before you move on}]
What does defekt mean? (\textquotedblleft{}Faulty, not working (of a device)\textquotedblright{}.) Say it once more.
\end{tcolorbox}
\section[(dialogue)]{Words from chapters 422–429 — a review}

\label{lesson:GE-R429-review-422-429}



\begin{quote}
Say the German for the cruise and to dock.
\end{quote}

\subsection*{Your turn}
Say these aloud:

\begin{itemize}
\raggedright
  \item the cruise, to dock, to cast off, to rock, the captain
  \item the anchor, to anchor, the lighthouse, the sailor, to capsize
  \item the inn, the holiday flat, to stay the night, the overnight stay, the half board
  \item the campsite, the caravan, the camper van, to put up, to take down
  \item the reception, to arrive, to leave, the stay, to check out
  \item the air conditioning, the heating, the radiator, to turn up, to shut off
  \item the drain, to block, to drip, leaky, the plumber
  \item the mould, to flicker, to jam, to squeak, faulty
\end{itemize}

\begin{tcolorbox}[breakable,colback=teal!4,colframe=teal!35!black,title={Before you move on}]
The cruise? (\textbf{die Kreuzfahrt}.) To dock? (\textbf{anlegen}.) To cast off? (\textbf{ablegen}.) To rock? (\textbf{schaukeln}.) The captain? (\textbf{der Kapitän}.) The anchor? (\textbf{der Anker}.) To anchor? (\textbf{ankern}.) The lighthouse? (\textbf{der Leuchtturm}.) The sailor? (\textbf{der Matrose}.) To capsize? (\textbf{kentern}.) The inn? (\textbf{der Gasthof}.) The holiday flat? (\textbf{die Ferienwohnung}.) To stay the night? (\textbf{übernachten}.) The overnight stay? (\textbf{die Übernachtung}.) The half board? (\textbf{die Halbpension}.) The campsite? (\textbf{der Campingplatz}.) The caravan? (\textbf{der Wohnwagen}.) The camper van? (\textbf{das Wohnmobil}.) To put up? (\textbf{aufbauen}.) To take down? (\textbf{abbauen}.) The reception? (\textbf{die Rezeption}.) To arrive? (\textbf{anreisen}.) To leave? (\textbf{abreisen}.) The stay? (\textbf{der Aufenthalt}.) To check out? (\textbf{auschecken}.) The air conditioning? (\textbf{die Klimaanlage}.) The heating? (\textbf{die Heizung}.) The radiator? (\textbf{der Heizkörper}.) To turn up? (\textbf{aufdrehen}.) To shut off? (\textbf{zudrehen}.) The drain? (\textbf{der Abfluss}.) To block? (\textbf{verstopfen}.) To drip? (\textbf{tropfen}.) Leaky? (\textbf{undicht}.) The plumber? (\textbf{der Klempner}.) The mould? (\textbf{der Schimmel}.) To flicker? (\textbf{flackern}.) To jam? (\textbf{klemmen}.) To squeak? (\textbf{quietschen}.) Faulty? (\textbf{defekt}.)
\end{tcolorbox}
